\subsection{Maximum on arbitrary connected algebraic curves}

\begin{lemma}[Closed rational sections of finite-type schemes]
\label{lem:fc-closed-rational-section}
\leanok
Every nonempty scheme locally of finite type over an algebraically closed
field has an original rational section that is a closed immersion, without
a separatedness hypothesis.
\end{lemma}
\begin{proof}
Choose the actual closed finite-extension field point constructed by an
affine maximal quotient and the Jacobson property. Its extension is the
original base field because that field is algebraically closed. Precompose
with the spectrum of the inverse base-field algebra isomorphism. This map
is a closed immersion since its ring map is surjective. Composing the two
closed immersions gives the required section and retains the exact original
structure triangle.
\end{proof}

\begin{lemma}[Integral reduced support closures retain the curve dimension]
\label{lem:fc-integral-closed-curve-support-lift}
\leanok
Every whole native point on a scheme of global dimension at most one lifts
to an actual integral closed subscheme of global dimension at most one.
\end{lemma}
\begin{proof}
Take the reduced vanishing-ideal subscheme of the singleton closure of its
original support. The singleton closure is nonempty and irreducible, so this
constructed closed subscheme is integral. A closed embedding does not increase
global dimension. The original support belongs to its own singleton closure;
the genuine closed-immersion native range theorem therefore lifts the whole
original native point to this subscheme.
\end{proof}

\begin{theorem}[Maximum on every original connected finite-type curve]
\label{thm:fc-connected-finite-type-curve-maximum}
\leanok
\uses{lem:fc-closed-rational-section,lem:fc-integral-closed-curve-support-lift,lem:fc-finite-type-closed-support-maximum}
Over a complete algebraically closed ultrametric normed field, a native locally
psh function on any connected finite-type scheme of global dimension at most
one is constant if it attains its global maximum. This includes every smooth
connected affine curve and also allows singular, reducible and nonreduced
curves and the trivial base norm.
\end{theorem}
\begin{proof}
Lift a maximizing point to the constructed integral reduced closure of its
support. The integral curve maximum theorem makes the pulled-back function
constant there. Choose a genuine closed rational section of this integral
finite-type subscheme. Its image in the original scheme has closed support
and the original maximum value.

For an arbitrary original native point construct its integral reduced support
closure and a closed rational section there. The image of this second section
has closed support in the original scheme, so connected finite-type closed
curve propagation gives it the same maximum value. The pulled-back function
therefore attains its global maximum at this section. The integral curve
maximum theorem equates this value with that of the lifted arbitrary point.
The original closed-immersion projection recovers the entire original point,
which proves the assertion without a smoothness or integrality input.
\end{proof}
